\documentclass[a5paper,10pt]{article}

% https://github.com/InDevRus/filippov-solutions

% Нумерация страниц
\pagenumbering{gobble}

% Настройка полей
\usepackage{geometry}
\geometry{left=1cm}
\geometry{right=1cm}
\geometry{top=1cm}
\geometry{bottom=1.5cm}

% Вставка таблицы
\usepackage{array}
\usepackage{tabularx}

% Инструменты выравнивания формул
\usepackage{amsmath}
\usepackage{mathtools}

% Диакритика в математике
\usepackage{amsfonts}

% Вычисления размеров на ходу
\usepackage{calc}

% Удобные записи производных
\usepackage[italic=true]{derivative}

% Настройки сносок
\usepackage{enotez}

% Длинные знаки векторов
\usepackage{anyfontsize}
\usepackage[b]{esvect}

% Вставка картинок
\usepackage{graphicx}

% Вставка красной строки
\usepackage{indentfirst}
\setlength{\parindent}{1em}

% Условия в командах
\usepackage{xifthen}

% Луа-скрипты
\usepackage{luacode}

% Настройка команд отдельно в математическом и текстовом режимах
\usepackage{mathcommand}

% Для повторения LaTeX кода
\usepackage{multido}

% Конфигурация отступов формул
\delimitershortfall-1sp
\usepackage{mleftright}
\mleftright

% Растягивание объектов
\usepackage{relsize}

% Обтекание текстом
\usepackage{wrapfig}
\usepackage{subcaption}
\usepackage{floatflt}

% Поддержка ввода кириллицы
\usepackage[english, russian]{babel}

% Настройка корневой позиции для компилятора
\usepackage{import}
\providecommand*{\ProjectRootPath}{..}

\import{\ProjectRootPath/preambles/macros/}{theorems}

\import{\ProjectRootPath/preambles/fonts}{font_setup}

% Знак градуса
\usepackage{siunitx}

\import{\ProjectRootPath/preambles/macros/}{operators}
\import{\ProjectRootPath/preambles/macros/}{redefinitions}

\import{\ProjectRootPath/preambles/fonts}{letters_setup}
\import{\ProjectRootPath/preambles/fonts/adjustments}{subsuperscripts}

\import{\ProjectRootPath/preambles/custom}{frames}
\import{\ProjectRootPath/preambles/custom}{formatting}
\import{\ProjectRootPath/preambles/custom}{list_setup}

% TODO: перевести команды на современный стиль объявления

\begin{document}
    $\br{\y - z} \vparder{z}{x} + \br{z - x} \vparder{z}{y} = x - y$.
    $\dvfrac {\di x} {y - z}
    = \dvfrac {\di y} {z - x}
    = \dvfrac {\di z} {x - y}$.

    Затемим, что в сумме знаменатели этих дробей зануляются, 
    так что можно скомбинировать, например, так
    $\dvfrac{\di x + \di y + \di z + \di x}{y - z} = \dvfrac {\di x} {y - z}$,
    откуда уже выводим интеграл
    $x + y + z = A$.

    Следующий интеграл неочевиден, так что вычислим его прямым методом с учётом найденного выше
    $\di z \br{\y - z} = \di x \br{x - \y}$. 
    $\di z \br{A - x - 2z} = \di x \br{2x + z - A}$. Далее,
    $2z \di z + 2x \di x + \br{x \di z + z \di x} - A\br{\di x + \di z} = 0$.
    $\pow{x}{2} + \pow{z}{2} + xz - Ax - Az = B$.
    Подставив $A = x + y + z$ и упростив, получаем
    $xy + yz + xz = B$.
    
    Отдельно отметим, что вместо полученного интеграла можно в том числе взять
    $C = \pow{x}{2} + \pow{y}{2} + \pow{z}{2} = \pow{A}{2} - 2B$.

    Итак, $x + y + z = A$, $xy + yz + xz = B$. При $z = y = -x$ имеем
    $-x = A$ и $-\pow{x}{2} = B$, так что $B + \pow{A}{2} = 0$.
    $\br{x + y + z}^2 + xy + yz + xz = 0$, либо же
    $\pow{x}{2} + \pow{y}{2} + \pow{z}{2} = 3 \br{x + y + z}^2$.

\end{document}